\scanpage{186}
% archive foliation: 185
désignant la ponctuation. Plus généralement, si on a un autre foncteur
\inserted{co}simplicial
\diagram{\begin{tikzcd}
f^0\arrow[r,shift left=.5ex]\arrow[r,shift right=.5ex]&
f^1\arrow[r,shift left=1ex]\arrow[r]\arrow[r,shift right=1ex]&f^2
\end{tikzcd}}
\deleted{\unsure{on va de}} et un hom \inserted{co}simplicial $\unsure{\zeta|\xi}$
\diagram{$\unsure{\zeta|\xi}:$\quad
\begin{tikzcd}[row sep=large]
F^0\arrow[r,shift left=.5ex]\arrow[r,shift right=.5ex]\arrow[d,"\unsure{\zeta^0|\xi^0}"']&
F^1\arrow[r,shift left=1ex]\arrow[r]\arrow[r,shift right=1ex]\arrow[d,"\unsure{\zeta^1|\xi^1}"']&
F^2\arrow[d,"\unsure{\zeta^2|\xi^2}"']\\
f^0\arrow[r,shift left=.5ex]\arrow[r,shift right=.5ex]&
f^1\arrow[r,shift left=1ex]\arrow[r]\arrow[r,shift right=1ex]&f^2
\end{tikzcd}}
on définit $\Pi^0(F^*,\unsure{\zeta|\xi};E)$ et
$\Pi^1(F^*,\unsure{\zeta|\xi};E)$.

Si $F^0,F^1,F^2$ sont $T$-représentables, et dans le cas où on se donne
des $\unsure{\zeta|\xi}:F^*\to f^*$, si $f^*$ est $T$-représentable,
on ramène ceci à la donnée d'un \deleted{objet} \uline{espace simplicial}
\diagram{\begin{tikzcd}
K_0&K_1\arrow[l,shift left=.5ex]\arrow[l,shift right=.5ex]&
K_2\arrow[l,shift left=1ex]\arrow[l]\arrow[l,shift right=1ex]
\end{tikzcd}}
resp d'un \uline{espace simplicial} \uuline{ponctué} (par contraction)
\diagram{\begin{tikzcd}
\overline K_0&\overline K_1\arrow[l,shift left=.5ex]\arrow[l,shift right=.5ex]&
\overline K_2\arrow[l,shift left=1ex]\arrow[l]\arrow[l,shift right=1ex]
\end{tikzcd}}
par un procédé standard \unsure{----}.

Il y a lieu de généraliser ces constructions au cas où
\unsure{sur des} $K_i$, resp. des $\overline K_i$, on a des catégories
galoisiennes attachées à chaque pt,
\deleted{\unsure{(connexes, les $F^i$}}
pour décrire des relations de descente à la
\uline{Van Kampen} \unsure{-----}
